% comparison_FR.tex — WAC confronted with KAPE and Eric Zimmerman's tools, French edition.
% The English edition is comparison_EN.tex, in this folder: both are maintained together,
% section for section.
%
% Build (LuaLaTeX, as the editor's recipe; build artefacts removed after):
%     ../build-docs.sh comparison_FR.tex
% The cover page is ../cover_comparison_FR.png, derived from ../cover_FR.png by
% ../retitle-cover.py. The layout is ../wac-style.tex, shared with every WAC document.
%
% Every figure comes from backlog/B-015 and the runs of vmtest/ez-compare.sh and
% vmtest/kape-compare.sh, which reproduce them.

\documentclass[11pt,a4paper,oneside]{report}

% ---------------------------------------------------------------- language and labels
\usepackage[french]{babel}
\newcommand{\wacdoctitle}{Face à KAPE et aux outils EZ}
\newcommand{\wacpageof}{sur}
\newcommand{\wacnotetitle}{À retenir}
\newcommand{\wacwarningtitle}{Limite}
% Layout shared with every WAC document.
% wac-style.tex — layout shared by every WAC document (user manuals, testing
% guides), French and English editions alike: one copy, so that they cannot
% drift apart.
%
% A document sets, BEFORE % wac-style.tex — layout shared by every WAC document (user manuals, testing
% guides), French and English editions alike: one copy, so that they cannot
% drift apart.
%
% A document sets, BEFORE % wac-style.tex — layout shared by every WAC document (user manuals, testing
% guides), French and English editions alike: one copy, so that they cannot
% drift apart.
%
% A document sets, BEFORE \input{../wac-style.tex}:
%   \usepackage[<language>]{babel}
%   \newcommand{\wacdoctitle}{...}      running header: "WAC — <title>"
%   \newcommand{\wacpageof}{...}        footer: "Page N <of> M"
%   \newcommand{\wacnotetitle}{...}     default title of the "note" box
%   \newcommand{\wacwarningtitle}{...}  default title of the "attention" box
% and AFTER it its \hypersetup (PDF title, author, subject).

% ---------------------------------------------------------------- fonts, language
\usepackage{fontspec}
\setmainfont{Noto Sans}
\setmonofont{DejaVu Sans Mono}[Scale=0.84]

% ---------------------------------------------------------------- layout
\usepackage[top=26mm,bottom=26mm,left=24mm,right=22mm,headheight=14pt]{geometry}
\usepackage{graphicx}
\graphicspath{{../}}
\usepackage{xcolor}
\usepackage{tikz}
\usepackage{booktabs,tabularx,longtable,array}
\usepackage{enumitem}
\usepackage{fancyhdr}
\usepackage{titlesec}
\usepackage[most]{tcolorbox}
\usepackage{listings}
\usepackage{url}
\usepackage{lastpage}
\usepackage[hidelinks,unicode,pdfpagelabels]{hyperref}

% Colours sampled from the cover page
\definecolor{marine}{HTML}{0A1A33}
\definecolor{bleu}{HTML}{2F6DB5}
\definecolor{accent}{HTML}{1F66D1}
\definecolor{fond}{HTML}{EEF3FA}
\definecolor{alerte}{HTML}{B3261E}
\definecolor{fondalerte}{HTML}{FBEDEC}
\definecolor{gris}{HTML}{5B6573}

\setlength{\parindent}{0pt}
\setlength{\parskip}{0.55em}
\linespread{1.08}
\setlist{itemsep=0.25em,topsep=0.35em}
\urlstyle{tt}
\renewcommand{\arraystretch}{1.25}
% Left-aligned table columns: justified narrow columns opened wide gaps
% between words.
\newcolumntype{L}[1]{>{\raggedright\arraybackslash}p{#1}}
% Monospaced file and field names cannot be hyphenated: let the line's spaces
% stretch instead.
\setlength{\emergencystretch}{3em}

% Headings
\titleformat{\chapter}[display]
  {\color{marine}\bfseries\raggedright}
  {\color{bleu}\large\MakeUppercase{\chaptertitlename}~\thechapter}
  {0.3em}
  {\Huge}[{\color{accent}\titlerule[2pt]}]
\titlespacing*{\chapter}{0pt}{0pt}{1.6em}
\titleformat{\section}{\color{marine}\Large\bfseries}{\color{bleu}\thesection}{0.8em}{}
\titleformat{\subsection}{\color{marine}\large\bfseries}{\color{bleu}\thesubsection}{0.7em}{}
\titleformat{\subsubsection}{\color{bleu}\normalsize\bfseries}{}{0em}{}

% Headers and footers
\pagestyle{fancy}
\fancyhf{}
\fancyhead[L]{\small\color{gris}WAC — \wacdoctitle}
\fancyhead[R]{\small\color{gris}\leftmark}
% Footer: "Page N of M" (the word set by \wacpageof), readable. A bare, small, light-grey number went
% unnoticed.
\newcommand{\piedpage}{\small\color{marine}\textbf{Page \thepage}\ \wacpageof\ \pageref*{LastPage}}
\fancyfoot[C]{\piedpage}
\renewcommand{\footrulewidth}{0.4pt}
\renewcommand{\footrule}{\hbox to\headwidth{\color{fond}\leaders\hrule height \footrulewidth\hfill}}
\renewcommand{\headrulewidth}{0.4pt}
\renewcommand{\headrule}{\hbox to\headwidth{\color{fond}\leaders\hrule height \headrulewidth\hfill}}
\renewcommand{\chaptermark}[1]{\markboth{#1}{}}
\fancypagestyle{plain}{\fancyhf{}\fancyfoot[C]{\piedpage}\renewcommand{\headrulewidth}{0pt}}

% Boxes
\newtcolorbox{note}[1][\wacnotetitle]{enhanced,breakable,colback=fond,colframe=bleu,
  boxrule=0pt,leftrule=3pt,arc=0pt,left=8pt,right=8pt,top=6pt,bottom=6pt,
  title={\bfseries #1},coltitle=bleu,colbacktitle=fond,attach title to upper={\par\smallskip}}
\newtcolorbox{attention}[1][\wacwarningtitle]{enhanced,breakable,colback=fondalerte,colframe=alerte,
  boxrule=0pt,leftrule=3pt,arc=0pt,left=8pt,right=8pt,top=6pt,bottom=6pt,
  title={\bfseries #1},coltitle=alerte,colbacktitle=fondalerte,attach title to upper={\par\smallskip}}

% Code
\lstset{basicstyle=\ttfamily\small,columns=fullflexible,keepspaces=true,
  backgroundcolor=\color{fond},frame=l,framerule=3pt,rulecolor=\color{bleu},
  xleftmargin=9pt,framexleftmargin=6pt,aboveskip=0.8em,belowskip=0.8em,
  breaklines=true,breakatwhitespace=false,showstringspaces=false,
  literate={é}{{\'e}}1 {è}{{\`e}}1 {à}{{\`a}}1 {ê}{{\^e}}1 {ô}{{\^o}}1 {ç}{{\c c}}1
           {É}{{\'E}}1 {—}{{---}}1 {’}{{'}}1 {…}{{...}}1 {«}{{\guillemotleft}}1 {»}{{\guillemotright}}1}

% Breakable monospace. A path or a field name in \texttt is one word that no
% hyphenation pattern cuts: in a table cell it ran 190 pt past the column
% (ProgramData\Microsoft\NetFramework\BreadcrumbStore). As the \seqsplit of
% the EventsManager theme's tables, a line may break inside it, without a
% hyphen that would read as part of the name: after \ / . _ - : and between a
% lowercase letter or digit and a capital (AmcacheSha1|Covers|Bytes). Only
% inside \texttt — the classes are set once, the breaks switched on there.
\newXeTeXintercharclass\wacBreakAfter
\newXeTeXintercharclass\wacUpper
\newXeTeXintercharclass\wacLower
\XeTeXcharclass`\\=\wacBreakAfter
\XeTeXcharclass`\/=\wacBreakAfter
\XeTeXcharclass`\.=\wacBreakAfter
\XeTeXcharclass`\_=\wacBreakAfter
\XeTeXcharclass`\-=\wacBreakAfter
\XeTeXcharclass`\:=\wacBreakAfter
\count255=`A \loop\XeTeXcharclass\count255=\wacUpper\ifnum\count255<`Z\advance\count255 by 1\repeat
\count255=`a \loop\XeTeXcharclass\count255=\wacLower\ifnum\count255<`z\advance\count255 by 1\repeat
\count255=`0 \loop\XeTeXcharclass\count255=\wacLower\ifnum\count255<`9\advance\count255 by 1\repeat
\XeTeXinterchartoks\wacBreakAfter\wacUpper={\allowbreak}
\XeTeXinterchartoks\wacBreakAfter\wacLower={\allowbreak}
\XeTeXinterchartoks\wacBreakAfter 0={\allowbreak}
\XeTeXinterchartoks\wacLower\wacUpper={\allowbreak}
% \textbackslash is no character token: it carries its own break.
\newcommand{\wacbackslash}{\char92\allowbreak}
\let\wacplaintexttt\texttt
\renewcommand{\texttt}[1]{\wacplaintexttt{\XeTeXinterchartokenstate=1\let\textbackslash\wacbackslash#1}}

% Shortcuts
\newcommand{\wac}{\textsc{wac}}
\newcommand{\opt}[1]{\texttt{-{}-#1}}
\newcommand{\champ}[1]{\texttt{#1}}
\newcommand{\fichier}[1]{\texttt{#1}}

% ============================================================ COVER PAGE
% The cover images are 2:3 (1024 x 1536), taller than an A4 page once widened
% to it: centred, they lost 9 mm at the top — the CEC logo — and 9 mm at the
% bottom. Raised 4 mm above the top of the page instead: the logo, which starts
% 6 mm below the image's edge, stays whole, and the bottom row of text keeps
% 7 mm to the page's edge; only a decorative corner below it is cut.
% Usage: \coverpage{cover_FR.png}
\newcommand{\coverpage}[1]{%
  \begin{titlepage}
  \thispagestyle{empty}
  \begin{tikzpicture}[remember picture,overlay]
    \node[inner sep=0pt,anchor=north,yshift=4mm] at (current page.north)
      {\includegraphics[width=\paperwidth]{#1}};
  \end{tikzpicture}
  \end{titlepage}}
:
%   \usepackage[<language>]{babel}
%   \newcommand{\wacdoctitle}{...}      running header: "WAC — <title>"
%   \newcommand{\wacpageof}{...}        footer: "Page N <of> M"
%   \newcommand{\wacnotetitle}{...}     default title of the "note" box
%   \newcommand{\wacwarningtitle}{...}  default title of the "attention" box
% and AFTER it its \hypersetup (PDF title, author, subject).

% ---------------------------------------------------------------- fonts, language
\usepackage{fontspec}
\setmainfont{Noto Sans}
\setmonofont{DejaVu Sans Mono}[Scale=0.84]

% ---------------------------------------------------------------- layout
\usepackage[top=26mm,bottom=26mm,left=24mm,right=22mm,headheight=14pt]{geometry}
\usepackage{graphicx}
\graphicspath{{../}}
\usepackage{xcolor}
\usepackage{tikz}
\usepackage{booktabs,tabularx,longtable,array}
\usepackage{enumitem}
\usepackage{fancyhdr}
\usepackage{titlesec}
\usepackage[most]{tcolorbox}
\usepackage{listings}
\usepackage{url}
\usepackage{lastpage}
\usepackage[hidelinks,unicode,pdfpagelabels]{hyperref}

% Colours sampled from the cover page
\definecolor{marine}{HTML}{0A1A33}
\definecolor{bleu}{HTML}{2F6DB5}
\definecolor{accent}{HTML}{1F66D1}
\definecolor{fond}{HTML}{EEF3FA}
\definecolor{alerte}{HTML}{B3261E}
\definecolor{fondalerte}{HTML}{FBEDEC}
\definecolor{gris}{HTML}{5B6573}

\setlength{\parindent}{0pt}
\setlength{\parskip}{0.55em}
\linespread{1.08}
\setlist{itemsep=0.25em,topsep=0.35em}
\urlstyle{tt}
\renewcommand{\arraystretch}{1.25}
% Left-aligned table columns: justified narrow columns opened wide gaps
% between words.
\newcolumntype{L}[1]{>{\raggedright\arraybackslash}p{#1}}
% Monospaced file and field names cannot be hyphenated: let the line's spaces
% stretch instead.
\setlength{\emergencystretch}{3em}

% Headings
\titleformat{\chapter}[display]
  {\color{marine}\bfseries\raggedright}
  {\color{bleu}\large\MakeUppercase{\chaptertitlename}~\thechapter}
  {0.3em}
  {\Huge}[{\color{accent}\titlerule[2pt]}]
\titlespacing*{\chapter}{0pt}{0pt}{1.6em}
\titleformat{\section}{\color{marine}\Large\bfseries}{\color{bleu}\thesection}{0.8em}{}
\titleformat{\subsection}{\color{marine}\large\bfseries}{\color{bleu}\thesubsection}{0.7em}{}
\titleformat{\subsubsection}{\color{bleu}\normalsize\bfseries}{}{0em}{}

% Headers and footers
\pagestyle{fancy}
\fancyhf{}
\fancyhead[L]{\small\color{gris}WAC — \wacdoctitle}
\fancyhead[R]{\small\color{gris}\leftmark}
% Footer: "Page N of M" (the word set by \wacpageof), readable. A bare, small, light-grey number went
% unnoticed.
\newcommand{\piedpage}{\small\color{marine}\textbf{Page \thepage}\ \wacpageof\ \pageref*{LastPage}}
\fancyfoot[C]{\piedpage}
\renewcommand{\footrulewidth}{0.4pt}
\renewcommand{\footrule}{\hbox to\headwidth{\color{fond}\leaders\hrule height \footrulewidth\hfill}}
\renewcommand{\headrulewidth}{0.4pt}
\renewcommand{\headrule}{\hbox to\headwidth{\color{fond}\leaders\hrule height \headrulewidth\hfill}}
\renewcommand{\chaptermark}[1]{\markboth{#1}{}}
\fancypagestyle{plain}{\fancyhf{}\fancyfoot[C]{\piedpage}\renewcommand{\headrulewidth}{0pt}}

% Boxes
\newtcolorbox{note}[1][\wacnotetitle]{enhanced,breakable,colback=fond,colframe=bleu,
  boxrule=0pt,leftrule=3pt,arc=0pt,left=8pt,right=8pt,top=6pt,bottom=6pt,
  title={\bfseries #1},coltitle=bleu,colbacktitle=fond,attach title to upper={\par\smallskip}}
\newtcolorbox{attention}[1][\wacwarningtitle]{enhanced,breakable,colback=fondalerte,colframe=alerte,
  boxrule=0pt,leftrule=3pt,arc=0pt,left=8pt,right=8pt,top=6pt,bottom=6pt,
  title={\bfseries #1},coltitle=alerte,colbacktitle=fondalerte,attach title to upper={\par\smallskip}}

% Code
\lstset{basicstyle=\ttfamily\small,columns=fullflexible,keepspaces=true,
  backgroundcolor=\color{fond},frame=l,framerule=3pt,rulecolor=\color{bleu},
  xleftmargin=9pt,framexleftmargin=6pt,aboveskip=0.8em,belowskip=0.8em,
  breaklines=true,breakatwhitespace=false,showstringspaces=false,
  literate={é}{{\'e}}1 {è}{{\`e}}1 {à}{{\`a}}1 {ê}{{\^e}}1 {ô}{{\^o}}1 {ç}{{\c c}}1
           {É}{{\'E}}1 {—}{{---}}1 {’}{{'}}1 {…}{{...}}1 {«}{{\guillemotleft}}1 {»}{{\guillemotright}}1}

% Breakable monospace. A path or a field name in \texttt is one word that no
% hyphenation pattern cuts: in a table cell it ran 190 pt past the column
% (ProgramData\Microsoft\NetFramework\BreadcrumbStore). As the \seqsplit of
% the EventsManager theme's tables, a line may break inside it, without a
% hyphen that would read as part of the name: after \ / . _ - : and between a
% lowercase letter or digit and a capital (AmcacheSha1|Covers|Bytes). Only
% inside \texttt — the classes are set once, the breaks switched on there.
\newXeTeXintercharclass\wacBreakAfter
\newXeTeXintercharclass\wacUpper
\newXeTeXintercharclass\wacLower
\XeTeXcharclass`\\=\wacBreakAfter
\XeTeXcharclass`\/=\wacBreakAfter
\XeTeXcharclass`\.=\wacBreakAfter
\XeTeXcharclass`\_=\wacBreakAfter
\XeTeXcharclass`\-=\wacBreakAfter
\XeTeXcharclass`\:=\wacBreakAfter
\count255=`A \loop\XeTeXcharclass\count255=\wacUpper\ifnum\count255<`Z\advance\count255 by 1\repeat
\count255=`a \loop\XeTeXcharclass\count255=\wacLower\ifnum\count255<`z\advance\count255 by 1\repeat
\count255=`0 \loop\XeTeXcharclass\count255=\wacLower\ifnum\count255<`9\advance\count255 by 1\repeat
\XeTeXinterchartoks\wacBreakAfter\wacUpper={\allowbreak}
\XeTeXinterchartoks\wacBreakAfter\wacLower={\allowbreak}
\XeTeXinterchartoks\wacBreakAfter 0={\allowbreak}
\XeTeXinterchartoks\wacLower\wacUpper={\allowbreak}
% \textbackslash is no character token: it carries its own break.
\newcommand{\wacbackslash}{\char92\allowbreak}
\let\wacplaintexttt\texttt
\renewcommand{\texttt}[1]{\wacplaintexttt{\XeTeXinterchartokenstate=1\let\textbackslash\wacbackslash#1}}

% Shortcuts
\newcommand{\wac}{\textsc{wac}}
\newcommand{\opt}[1]{\texttt{-{}-#1}}
\newcommand{\champ}[1]{\texttt{#1}}
\newcommand{\fichier}[1]{\texttt{#1}}

% ============================================================ COVER PAGE
% The cover images are 2:3 (1024 x 1536), taller than an A4 page once widened
% to it: centred, they lost 9 mm at the top — the CEC logo — and 9 mm at the
% bottom. Raised 4 mm above the top of the page instead: the logo, which starts
% 6 mm below the image's edge, stays whole, and the bottom row of text keeps
% 7 mm to the page's edge; only a decorative corner below it is cut.
% Usage: \coverpage{cover_FR.png}
\newcommand{\coverpage}[1]{%
  \begin{titlepage}
  \thispagestyle{empty}
  \begin{tikzpicture}[remember picture,overlay]
    \node[inner sep=0pt,anchor=north,yshift=4mm] at (current page.north)
      {\includegraphics[width=\paperwidth]{#1}};
  \end{tikzpicture}
  \end{titlepage}}
:
%   \usepackage[<language>]{babel}
%   \newcommand{\wacdoctitle}{...}      running header: "WAC — <title>"
%   \newcommand{\wacpageof}{...}        footer: "Page N <of> M"
%   \newcommand{\wacnotetitle}{...}     default title of the "note" box
%   \newcommand{\wacwarningtitle}{...}  default title of the "attention" box
% and AFTER it its \hypersetup (PDF title, author, subject).

% ---------------------------------------------------------------- fonts, language
\usepackage{fontspec}
\setmainfont{Noto Sans}
\setmonofont{DejaVu Sans Mono}[Scale=0.84]

% ---------------------------------------------------------------- layout
\usepackage[top=26mm,bottom=26mm,left=24mm,right=22mm,headheight=14pt]{geometry}
\usepackage{graphicx}
\graphicspath{{../}}
\usepackage{xcolor}
\usepackage{tikz}
\usepackage{booktabs,tabularx,longtable,array}
\usepackage{enumitem}
\usepackage{fancyhdr}
\usepackage{titlesec}
\usepackage[most]{tcolorbox}
\usepackage{listings}
\usepackage{url}
\usepackage{lastpage}
\usepackage[hidelinks,unicode,pdfpagelabels]{hyperref}

% Colours sampled from the cover page
\definecolor{marine}{HTML}{0A1A33}
\definecolor{bleu}{HTML}{2F6DB5}
\definecolor{accent}{HTML}{1F66D1}
\definecolor{fond}{HTML}{EEF3FA}
\definecolor{alerte}{HTML}{B3261E}
\definecolor{fondalerte}{HTML}{FBEDEC}
\definecolor{gris}{HTML}{5B6573}

\setlength{\parindent}{0pt}
\setlength{\parskip}{0.55em}
\linespread{1.08}
\setlist{itemsep=0.25em,topsep=0.35em}
\urlstyle{tt}
\renewcommand{\arraystretch}{1.25}
% Left-aligned table columns: justified narrow columns opened wide gaps
% between words.
\newcolumntype{L}[1]{>{\raggedright\arraybackslash}p{#1}}
% Monospaced file and field names cannot be hyphenated: let the line's spaces
% stretch instead.
\setlength{\emergencystretch}{3em}

% Headings
\titleformat{\chapter}[display]
  {\color{marine}\bfseries\raggedright}
  {\color{bleu}\large\MakeUppercase{\chaptertitlename}~\thechapter}
  {0.3em}
  {\Huge}[{\color{accent}\titlerule[2pt]}]
\titlespacing*{\chapter}{0pt}{0pt}{1.6em}
\titleformat{\section}{\color{marine}\Large\bfseries}{\color{bleu}\thesection}{0.8em}{}
\titleformat{\subsection}{\color{marine}\large\bfseries}{\color{bleu}\thesubsection}{0.7em}{}
\titleformat{\subsubsection}{\color{bleu}\normalsize\bfseries}{}{0em}{}

% Headers and footers
\pagestyle{fancy}
\fancyhf{}
\fancyhead[L]{\small\color{gris}WAC — \wacdoctitle}
\fancyhead[R]{\small\color{gris}\leftmark}
% Footer: "Page N of M" (the word set by \wacpageof), readable. A bare, small, light-grey number went
% unnoticed.
\newcommand{\piedpage}{\small\color{marine}\textbf{Page \thepage}\ \wacpageof\ \pageref*{LastPage}}
\fancyfoot[C]{\piedpage}
\renewcommand{\footrulewidth}{0.4pt}
\renewcommand{\footrule}{\hbox to\headwidth{\color{fond}\leaders\hrule height \footrulewidth\hfill}}
\renewcommand{\headrulewidth}{0.4pt}
\renewcommand{\headrule}{\hbox to\headwidth{\color{fond}\leaders\hrule height \headrulewidth\hfill}}
\renewcommand{\chaptermark}[1]{\markboth{#1}{}}
\fancypagestyle{plain}{\fancyhf{}\fancyfoot[C]{\piedpage}\renewcommand{\headrulewidth}{0pt}}

% Boxes
\newtcolorbox{note}[1][\wacnotetitle]{enhanced,breakable,colback=fond,colframe=bleu,
  boxrule=0pt,leftrule=3pt,arc=0pt,left=8pt,right=8pt,top=6pt,bottom=6pt,
  title={\bfseries #1},coltitle=bleu,colbacktitle=fond,attach title to upper={\par\smallskip}}
\newtcolorbox{attention}[1][\wacwarningtitle]{enhanced,breakable,colback=fondalerte,colframe=alerte,
  boxrule=0pt,leftrule=3pt,arc=0pt,left=8pt,right=8pt,top=6pt,bottom=6pt,
  title={\bfseries #1},coltitle=alerte,colbacktitle=fondalerte,attach title to upper={\par\smallskip}}

% Code
\lstset{basicstyle=\ttfamily\small,columns=fullflexible,keepspaces=true,
  backgroundcolor=\color{fond},frame=l,framerule=3pt,rulecolor=\color{bleu},
  xleftmargin=9pt,framexleftmargin=6pt,aboveskip=0.8em,belowskip=0.8em,
  breaklines=true,breakatwhitespace=false,showstringspaces=false,
  literate={é}{{\'e}}1 {è}{{\`e}}1 {à}{{\`a}}1 {ê}{{\^e}}1 {ô}{{\^o}}1 {ç}{{\c c}}1
           {É}{{\'E}}1 {—}{{---}}1 {’}{{'}}1 {…}{{...}}1 {«}{{\guillemotleft}}1 {»}{{\guillemotright}}1}

% Breakable monospace. A path or a field name in \texttt is one word that no
% hyphenation pattern cuts: in a table cell it ran 190 pt past the column
% (ProgramData\Microsoft\NetFramework\BreadcrumbStore). As the \seqsplit of
% the EventsManager theme's tables, a line may break inside it, without a
% hyphen that would read as part of the name: after \ / . _ - : and between a
% lowercase letter or digit and a capital (AmcacheSha1|Covers|Bytes). Only
% inside \texttt — the classes are set once, the breaks switched on there.
\newXeTeXintercharclass\wacBreakAfter
\newXeTeXintercharclass\wacUpper
\newXeTeXintercharclass\wacLower
\XeTeXcharclass`\\=\wacBreakAfter
\XeTeXcharclass`\/=\wacBreakAfter
\XeTeXcharclass`\.=\wacBreakAfter
\XeTeXcharclass`\_=\wacBreakAfter
\XeTeXcharclass`\-=\wacBreakAfter
\XeTeXcharclass`\:=\wacBreakAfter
\count255=`A \loop\XeTeXcharclass\count255=\wacUpper\ifnum\count255<`Z\advance\count255 by 1\repeat
\count255=`a \loop\XeTeXcharclass\count255=\wacLower\ifnum\count255<`z\advance\count255 by 1\repeat
\count255=`0 \loop\XeTeXcharclass\count255=\wacLower\ifnum\count255<`9\advance\count255 by 1\repeat
\XeTeXinterchartoks\wacBreakAfter\wacUpper={\allowbreak}
\XeTeXinterchartoks\wacBreakAfter\wacLower={\allowbreak}
\XeTeXinterchartoks\wacBreakAfter 0={\allowbreak}
\XeTeXinterchartoks\wacLower\wacUpper={\allowbreak}
% \textbackslash is no character token: it carries its own break.
\newcommand{\wacbackslash}{\char92\allowbreak}
\let\wacplaintexttt\texttt
\renewcommand{\texttt}[1]{\wacplaintexttt{\XeTeXinterchartokenstate=1\let\textbackslash\wacbackslash#1}}

% Shortcuts
\newcommand{\wac}{\textsc{wac}}
\newcommand{\opt}[1]{\texttt{-{}-#1}}
\newcommand{\champ}[1]{\texttt{#1}}
\newcommand{\fichier}[1]{\texttt{#1}}

% ============================================================ COVER PAGE
% The cover images are 2:3 (1024 x 1536), taller than an A4 page once widened
% to it: centred, they lost 9 mm at the top — the CEC logo — and 9 mm at the
% bottom. Raised 4 mm above the top of the page instead: the logo, which starts
% 6 mm below the image's edge, stays whole, and the bottom row of text keeps
% 7 mm to the page's edge; only a decorative corner below it is cut.
% Usage: \coverpage{cover_FR.png}
\newcommand{\coverpage}[1]{%
  \begin{titlepage}
  \thispagestyle{empty}
  \begin{tikzpicture}[remember picture,overlay]
    \node[inner sep=0pt,anchor=north,yshift=4mm] at (current page.north)
      {\includegraphics[width=\paperwidth]{#1}};
  \end{tikzpicture}
  \end{titlepage}}


\hypersetup{pdftitle={WAC — Face à KAPE et aux outils EZ},
  pdfauthor={Centre d'Excellence Cyberdéfense Aérospatiale (CEC)},
  pdfsubject={WAC confronté à KAPE et aux outils d'Eric Zimmerman : temps, précision, formats}}

\begin{document}

% ============================================================ COVER PAGE
\pagenumbering{roman}
\coverpage{cover_comparison_FR.png}

% ============================================================ INFORMATION PAGE
\thispagestyle{empty}
\vspace*{\fill}
{\color{marine}\Large\bfseries WAC — Windows Artefacts Collector}\par
{\color{bleu}Face à KAPE et aux outils d'Eric Zimmerman}\par
\medskip
\begin{tabular}{@{}ll@{}}
Version du logiciel & 3.0.0 \\
Édition du document & octobre 2026 \\
Éditeur & Centre d'Excellence Cyberdéfense Aérospatiale (CEC) \\
        & École de l'air et de l'espace, base aérienne 701, Salon-de-Provence \\
\end{tabular}

\bigskip
{\small\color{gris}Ce document confronte \wac{} aux outils de référence de la
communauté : KAPE pour la collecte, les outils en ligne de commande d'Eric
Zimmerman pour le décodage. Sur les mêmes pièces d'une même machine, il compare
les temps d'exécution, ce que chacun prélève, la justesse des valeurs --- chaque
écart étant tranché par une référence indépendante ---, les formats de sortie, et
les traces que chacun laisse sur la machine examinée. Il dit ce que la
confrontation a appris à \wac{}, et ce qu'elle lui reproche encore.}
\vspace*{2cm}
\clearpage

\pagenumbering{arabic}
\tableofcontents

% ============================================================ 1
\chapter{La démarche}

\section{Ce qui est comparé}

Deux questions, deux comparaisons.

\begin{description}[style=nextline,leftmargin=1em]
  \item[Collecte : \wac{} face à KAPE] Les deux outils collectent la même
        machine l'un après l'autre ; on compare ce que chacun copie, le temps
        qu'il y met, et l'empreinte SHA-1 des fichiers qu'ils copient tous deux.
  \item[Traces laissées : \wac{}, KAPE et les outils EZ] Chacun est lancé depuis
        le même état d'une même machine, sous Sysmon ; on relève ce qu'il y
        écrit, ce que Windows y écrit à cause de lui, et ce qu'un outil de
        surveillance en voit (chapitre~\ref{chap:traces}).
  \item[Décodage : \wac{} face aux outils EZ] Sur les \emph{mêmes pièces} --- la
        consigne d'une collecte \wac{} ---, \wac{} convertit la collecte entière
        et l'outil d'Eric Zimmerman correspondant décode, sur une copie, les
        pièces de son artefact ; les deux sorties sont comparées enregistrement
        par enregistrement et champ par champ.
\end{description}

\section{Qui a raison quand ils diffèrent}

Un écart entre deux outils ne dit pas lequel se trompe. Chaque écart a donc été
tranché par une troisième source, indépendante des deux :

\begin{itemize}[nosep]
  \item Windows lui-même : le XML qu'il rend pour un événement
        (\texttt{Get-WinEvent}), la date d'un fichier, l'identifiant d'objet
        NTFS (\texttt{fsutil objectid}), le registre en ligne, l'empreinte d'un
        fichier (\texttt{Get-FileHash}) ;
  \item une autre implémentation : \texttt{dissect.esedb} (Fox-IT) pour les bases
        ESE ;
  \item le format lui-même : numéros de séquence et sommes Marvin32 des journaux
        de ruche, lus octet par octet ;
  \item le recoupement de deux artefacts de la même machine : le Prefetch pour
        l'exécution, l'horodatage d'une ligne pour ses autres dates.
\end{itemize}

Chaque écart est ensuite classé : \emph{défaut de \wac{}} (corrigé), \emph{erreur
de l'outil EZ} (vérifiée), ou \emph{convention} de l'un des deux --- deux formes
d'une même valeur.

\section{Le banc d'essai}

\begin{tabular}{@{}ll@{}}
Machine & Windows 11 Professionnel (build 26200), machine virtuelle \\
Outils EZ & version 2026.5.0 (.NET 9.0.20) \\
KAPE & version 1.3.0.2, cible \texttt{KapeTriage} \\
\wac{} & version 3.0.0 et ses corrections issues de cette confrontation \\
Mesures & 29 et 30 octobre 2026 \\
Mesure des traces & clone de la machine, sans réseau, figé ; Sysmon 15.22 ; \\
                  & outils sur une clé USB (K:), une exécution par outil \\
\end{tabular}

Tout se rejoue : \fichier{vmtest/ez-compare.sh <artefact>} chronomètre l'outil
EZ sur une copie des pièces et compare sa sortie à la conversion de \wac{} ;
\fichier{vmtest/lib/timing-table.py} met les temps de chaque étape en regard ;
\fichier{vmtest/kape-compare.sh} fait de même pour la collecte ;
\fichier{vmtest/footprint/run-footprint.sh} rejoue la mesure des traces et
\fichier{footprint-diff.py} l'analyse. Les outils EZ ont
été lancés avec leur pleine précision de date (\opt{mp}, ou \opt{dt} au format
à sept décimales quand \opt{mp} ne suffit pas), pour ne pas leur imputer une
perte de précision qu'un réglage évite.

% ============================================================ 2
\chapter{Synthèse}

Pour chaque artefact, les sorties de \wac{} et de l'outil EZ sont appariées
enregistrement par enregistrement, puis comparées champ par champ. La dernière
colonne donne chaque écart trouvé et qui a raison, établi par une source
indépendante (chapitre~\ref{chap:precision}) ; un champ qui n'y figure pas a
la même valeur des deux côtés.

\begin{longtable}{@{}L{3.0cm}L{2.3cm}L{4.6cm}L{5.0cm}@{}}
\toprule
\textbf{Artefact --- outil} & \textbf{Comparés} & \textbf{Champs comparés} & \textbf{Écarts, et qui a raison} \\
\midrule\endhead
\$MFT --- MFTECmd & 358\,335 entrées & chemin, état (supprimé, dossier), 8 dates, taille, entrée MFT & tailles de 2\,010 fichiers compressés (WOF) et de 10 fichiers ouverts en écriture : \wac{}, comme Windows \\
EVTX --- EvtxECmd & 246\,096 événements & identifiant, date, fournisseur, machine, données de l'événement & 25 valeurs : \wac{}, comme Windows ; 835 événements Perflib qu'EvtxECmd ne lit pas ; journaux vidés qu'il ignore \\
Prefetch --- PECmd & 411 fichiers & exécutable, empreinte, exécutions et leurs dates, volumes, fichiers chargés & empreinte sans son zéro de tête chez PECmd (20) ; répertoires de deux volumes collés (1) ; date de volume nulle écrite 1601 (2) \\
Raccourcis --- LECmd & 67 fichiers & cible, dates au 100\,ns, taille, volume, machine, arguments & aucun \\
Jump lists --- JLECmd & 123 entrées & chemin, compteur, épinglage, dates, raccourci de chaque entrée & aucun \\
Shellbags --- SBECmd & 46 dossiers & nom, rang, dates de la clé et du dossier, entrée MFT & NodeSlot faux au premier niveau chez SBECmd (9) ; « exploré » toujours faux \\
Amcache --- AmcacheParser & 1\,890 entrées & les 124 colonnes communes des 8 tables & 5\,500 valeurs absentes de la ruche, inventées par AmcacheParser ; SHA-1 partiel non signalé \\
ShimCache --- AppCompat\-Cache\-Parser & 596 entrées & chemin, date, rang & aucun ; l'indicateur « exécuté » d'AppCompatCacheParser, trompeur, n'est pas relevé \\
Corbeille --- RBCmd & 130 fichiers & chemin d'origine, taille, date de suppression & aucun \\
SRUM --- SrumECmd & 115\,899 lignes & les 7 tables, toutes leurs colonnes & 8 lignes d'énergie décalées d'une ligne chez SrumECmd (dissect.esedb) \\
Registre --- RECmd & 9 artefacts & UserAssist, BAM, USBSTOR, comptes, Run, MuiCache, chemins tapés, recherches, documents de confiance & ruches ramenées à un état ancien par RECmd (BAM, UserAssist, USBSTOR, SAM, documents de confiance) \\
Collecte --- KAPE & 1\,430 fichiers pris des deux côtés & empreinte SHA-1 & 18 différents : fichiers que Windows a réécrits entre les deux collectes \\
\bottomrule
\end{longtable}

\begin{note}
Sur tout ce que les deux côtés lisent, \wac{} donne les mêmes valeurs que les
outils de référence, ou des valeurs plus justes, établies à chaque fois par une
source indépendante. La confrontation a aussi révélé 29 défauts de \wac{}, tous
corrigés (chapitre~\ref{chap:defauts}), et des lacunes de couverture,
toutes comblées depuis (section~\ref{sec:couverture}).

Lancé depuis une clé, \wac{} n'écrit rien sur la machine examinée : ni fichier,
ni valeur de registre, ni événement dans les journaux de Windows. Restent les
traces que Windows laisse de tout programme lancé (Amcache, ShimCache), et ce
qu'une surveillance en place enregistre de son exécution : Sysmon, l'audit de
création de processus, un EDR ; les ruches relues hors ligne le confirment.
KAPE, un programme .NET, écrit dans le registre de l'utilisateur, et ses
modules des clés dans SOFTWARE et des fichiers sur le volume système
(chapitre~\ref{chap:traces}).
\end{note}

% ============================================================ 3
\chapter{Temps d'exécution}\label{chap:temps}

\section{La méthode}

\wac{} ne se mesure pas artefact par artefact. Une conversion lancée pour un
seul artefact lui fait payer, à chaque fois, ce qu'il ne paie qu'une fois en
vrai : vérifier le sceau et les empreintes de la consigne, reconstruire la copie
de travail, réparer les ruches. Une collecte lancée pour un seul artefact lui
fait rouvrir et parcourir le volume à chaque fois. Les temps ci-dessous viennent
donc d'\emph{une} collecte complète et d'\emph{une} conversion complète.
\wac{} y note la durée de chacune de ses étapes, dans \fichier{investigation.json}
et \fichier{conversion.json}. Chaque artefact est mis en face de l'outil EZ qui
le décode. Les étapes que les autres outils n'ont pas sont données à part.
Le rapprochement est fait par \fichier{vmtest/lib/timing-table.py}.

\section{Collecte}

KAPE (cible \texttt{KapeTriage}) et \wac{} collectent la même machine, l'un
juste après l'autre.

\begin{longtable}{@{}L{6.2cm}L{2.2cm}L{2.6cm}L{4.0cm}@{}}
\toprule
\textbf{Outil, réglage} & \textbf{Durée} & \textbf{Fichiers} & \textbf{Volume} \\
\midrule\endhead
KAPE, \texttt{KapeTriage} & 53\,s & 1\,435 & 0,98\,Go \\
\wac{}, sans les exécutables (\texttt{binary: none}) & 226\,s & 16\,042 & 7,52\,Go \\
\quad dont les fichiers supprimés récupérés & 156\,s & 12\,529 & 4,29\,Go \\
\wac{}, kit livré (\texttt{binary: unverified}) & 626\,s & 72\,188 consignés & 8,50\,Go \\
\bottomrule
\end{longtable}

Les deux tiers de la collecte de \wac{} vont aux fichiers supprimés, que la
version 3.0.0 récupère par défaut (niveau~2, 4\,Gio copiés au plus) : KAPE
n'en récupère aucun. Sans eux, \wac{} prend en 64\,s ce que KAPE prend en
53\,s, et 2\,078 fichiers de plus (3\,513 en tout, 3,2\,Go). Ses étapes,
mesurées dans cette exécution :

\begin{longtable}{@{}L{8.4cm}L{2.0cm}L{4.6cm}@{}}
\toprule
\textbf{Étape de \wac{}} & \textbf{Durée} & \textbf{Chez KAPE} \\
\midrule\endhead
Contrôles : droits, place libre sur le support & 0,00\,s & aucun contrôle de place \\
Observation de la machine en marche : sessions, processus, réseau, services, horloge & 0,15\,s & --- \\
Ruches du système, en lecture brute NTFS ; le volume est ouvert une fois, et
  ses répertoires mis en cache pour toute la collecte & 2,34\,s & copie, lecture brute pour les fichiers verrouillés \\
Ruches des utilisateurs & 0,25\,s & idem \\
Fichiers des artefacts (journaux, Prefetch, raccourcis, navigateurs, traces
  ETW, courrier…) & 38,22\,s & idem \\
\$MFT, \$UsnJrnl, \$LogFile, \$Boot, \$Secure:\$SDS & 6,63\,s & idem \\
Secteurs et fichiers de démarrage (MBR, GPT, secteurs de démarrage, partition EFI) & 2,51\,s & --- \\
Recherche des volumes FAT et exFAT (clés USB) et lecture de leurs métadonnées & 0,04\,s & --- \\
Clichés instantanés, en lecture brute, sans le service VSS (aucun sur la machine) & 0,01\,s & option, par le service VSS \\
Fichiers de ressources des fournisseurs que citent les journaux, pour le
  texte des événements & 13,76\,s & --- \\
Fichiers supprimés : contenu des enregistrements libres de la \$MFT et noms
  restés dans les index & 156,10\,s & --- \\
Scellement : manifeste, empreinte SHA-256 du manifeste & 0,29\,s & journal de copie \\
Hors étapes : démarrage, écriture du journal d'investigation & 4,7\,s & \\
\bottomrule
\end{longtable}

Les trois empreintes de chaque pièce (MD5, SHA-1, SHA-256) et ses dates NTFS
sont prises pendant la lecture même : le scellement n'a plus qu'à écrire le
manifeste. Hors fichiers supprimés, les étapes propres à \wac{} --- contrôles,
observation de la machine, démarrage, volumes FAT, clichés instantanés,
ressources des fournisseurs, scellement --- coûtent 16,8\,s, dont 13,8\,s pour
les ressources des fournisseurs. Ce sont elles qui permettent, sur le poste
d'analyse, de rendre le texte de chaque événement comme Windows l'affiche. Le
niveau et le budget des fichiers supprimés se règlent dans \fichier{wac.yml}
(\texttt{deleted\_files}).

Avec les exécutables (réglage du kit livré), la collecte prend 626\,s, dont
334\,s pour lire chaque exécutable, bibliothèque, pilote et script du volume.
\wac{} en vérifie la signature Authenticode, en mémoire, contre son jeu de
confiance, et ne copie que ceux dont la signature ne tient pas. Les 40\,531
fichiers lus se répartissent ainsi :
\begin{itemize}[nosep]
  \item 25\,875 authentifiés comme Microsoft et 329 validés par une chaîne
        tierce : empreintes consignées, fichiers non copiés ;
  \item 11\,334 fichiers différentiels de Windows (MSDELTA), qui ne sont pas
        des exécutables : empreintes consignées, non copiés ;
  \item 2\,768 dont la signature ne tient pas, copiés (823\,Mio) ;
  \item 225 au contenu identique à une pièce déjà copiée : consignés,
        renvoyés à cette pièce, non copiés une seconde fois.
\end{itemize}
Les 41\,s de l'observation des processus
viennent, dans ce réglage, des empreintes de chaque module chargé, et les
6\,s du jeu de confiance de son chargement. KAPE ne fait aucun de ces
contrôles : il copie les fichiers de ses cibles sans les examiner.

\section{Conversion et décodage}

La conversion de \wac{} lit la consigne d'une collecte, sur le poste
d'analyse. Chaque outil EZ lit une copie des mêmes pièces : l'outil est
chronométré seul, la copie non. Cette copie n'est pas une précaution de forme.
Lancé sur la consigne elle-même, SrumECmd y a écrit (son \texttt{JetInit} pose
un point de contrôle et des journaux à côté de la base), et \wac{} a ensuite
refusé de convertir une consigne qui ne correspondait plus à son manifeste.

\begin{longtable}{@{}L{4.0cm}L{2.1cm}L{2.5cm}L{2.1cm}L{3.8cm}@{}}
\toprule
\textbf{Artefact} & \textbf{\wac{}} & \textbf{\wac{} \opt{binary}} & \textbf{Outil EZ} & \textbf{Outil} \\
\midrule\endhead
\$MFT (C: et G:) & 4,81\,s & 4,68\,s & 10,0\,s & MFTECmd (C: seul) \\
Journaux d'événements & 15,37\,s & 16,34\,s & 33,1\,s & EvtxECmd \\
Prefetch (411) & 0,36\,s & 2,78\,s & 3,0\,s & PECmd \\
Raccourcis & 0,04\,s & 0,19\,s & 4,4\,s & LECmd \\
Jump lists & 0,04\,s & 0,44\,s & 0,8\,s & JLECmd \\
Shellbags & 0,02\,s & 0,04\,s & 1,0\,s & SBECmd \\
ShimCache & 0,00\,s & 0,32\,s & 1,5\,s & AppCompat\-Cache\-Parser \\
Amcache & 0,12\,s & 11,14\,s & 0,8\,s & AmcacheParser \\
Corbeille & 0,02\,s & 0,07\,s & 0,6\,s & RBCmd \\
SRUM & 2,26\,s & 2,36\,s & 3,1\,s & SrumECmd \\
Registre (lot Kroll) & 0,49\,s & 2,45\,s & 15,9\,s & RECmd \\
\midrule
\textbf{Somme} & 23,5\,s & 40,8\,s & 74,2\,s & les 11 outils \\
\bottomrule
\end{longtable}

À ces étapes s'ajoutent celles que \wac{} paie une fois pour tous les artefacts,
et que les outils EZ n'ont pas :

\begin{longtable}{@{}L{11.0cm}L{4.0cm}@{}}
\toprule
\textbf{Étape commune de \wac{}} & \textbf{Durée} \\
\midrule\endhead
Vérification du sceau et des empreintes de chaque pièce & 37,2\,s \\
Reconstruction de la copie de travail depuis la consigne (5,9\,Go, dont
  4,3\,Go de fichiers supprimés), chaque copie vérifiée par son empreinte & 136,8\,s \\
Réparation des ruches de la copie de travail (rejeu de leurs journaux) & 0,5\,s \\
Fuseau, page de codes, profils et noms des comptes du suspect & 0,0\,s \\
\midrule
\textbf{Total} & 174,5\,s \\
\bottomrule
\end{longtable}

La conversion entière dure 209\,s : les 175\,s communs, les 23,5\,s des onze
artefacts comparés, et 10,8\,s pour les 45 autres, que les outils EZ ne lisent
pas (fichiers supprimés, journal des modifications, tâches planifiées, caches
de miniatures, connexions RDP, courrier, navigateurs…). Les onze outils EZ
prennent ensemble 74\,s, sans copie de travail ni vérification de la
consigne. Les fichiers supprimés pèsent surtout ici : leurs 4,3\,Go sont
vérifiés et recopiés comme toute pièce.

\begin{itemize}
  \item Une fois la copie de travail faite, chaque décodage de \wac{} est plus
        rapide que l'outil EZ correspondant : les ruches sont lues une fois,
        pour tous les artefacts du registre.
  \item L'étape la plus lourde est la copie de travail. Elle est la condition
        d'une analyse qui ne touche jamais la consigne, et la même copie est
        nécessaire aux outils EZ, comme SrumECmd l'a montré.
  \item Avec \opt{binary}, \wac{} calcule l'empreinte de chaque fichier que citent
        les artefacts : +2,4\,s pour le Prefetch (les bibliothèques chargées), +11,0\,s
        pour l'Amcache. Les outils EZ ne le font pas.
\end{itemize}

\section{Collecte depuis un même état}

Mesures de la comparaison des traces (chapitre~\ref{chap:traces}), les 29 et
30 octobre 2026 : chaque outil part du même état figé de la machine, lancé
depuis la clé ; la durée est celle de la tâche qui le lance, entre parenthèses
celle que l'outil annonce.

\begin{longtable}{@{}L{5.6cm}L{2.6cm}L{3.4cm}L{3.4cm}@{}}
\toprule
\textbf{Outil, réglage} & \textbf{Durée} & \textbf{Pris} & \textbf{Sur la clé} \\
\midrule\endhead
KAPE, \texttt{!SANS\_Triage} & 167\,s (136\,s) & 2\,674 fichiers & 2\,664 fichiers, 1,33\,Go \\
KAPE, \texttt{!SANS\_Triage} et modules \texttt{!EZParser} & 304\,s (289\,s) & 2\,673 fichiers, 14 outils & 2\,717 fichiers, 1,86\,Go \\
Outils EZ sur la machine vivante (6 outils) & 106\,s & \$MFT, journaux, ruches, Prefetch, Amcache, ShimCache & 36 CSV, 0,60\,Go \\
\wac{}, collecte (réglage par défaut) & 927\,s (924\,s) & 19\,350 pièces, 71\,249 consignées & 22\,004 fichiers, 9,84\,Go \\
\bottomrule
\end{longtable}

La collecte de \wac{} est la plus longue parce qu'elle lit \emph{chaque} exécutable
de la machine --- la conversion se fera ailleurs, et elle ne sait pas encore
lesquels les artefacts citeront --- et parce qu'elle récupère les fichiers
supprimés. Elle calcule les empreintes des exécutables, vérifie leur signature
et ne copie que ceux dont la signature ne tient pas ; d'où 71\,249 pièces
consignées pour 19\,350 copiées. La conversion, sur le poste d'analyse, ne
relit rien de la machine (chapitre~\ref{chap:traces}). KAPE copie les fichiers
de ses cibles sans les examiner.

% ============================================================ 4
\chapter{Ce que chacun prélève}\label{chap:prelevement}

\section{Les mêmes pièces, des fichiers différents}

Dernière comparaison, le 29 octobre 2026 : KAPE avec sa cible
\texttt{KapeTriage}, et \wac{} sans les exécutables (\texttt{binary: none}),
l'un juste après l'autre ; les fichiers sont appariés par leur chemin
d'origine (manifeste de \wac{}, journal de copie de KAPE). KAPE prend 1\,435
fichiers, \wac{} 16\,040 ; 1\,430 sont pris par les deux.

KAPE ne prend que 5 fichiers que \wac{} ne prend pas : des \fichier{desktop.ini},
qui ne fixent que l'icône et le nom affiché d'un dossier. \wac{} prend en plus
14\,610 fichiers : les fichiers supprimés qu'il récupère (12\,529) ; les journaux
et réglages de OneDrive (811) ; les fichiers de ressources des fournisseurs
d'événements (285), qui donnent le texte des événements ; les traces ETW
(227) ; Defender (186) ; les secteurs et fichiers de démarrage (158) ; le
courrier de Thunderbird (156) ; les raccourcis hors de Recent (76) ; l'index
de recherche de Windows et les caches de miniatures (53 en tout) ; d'autres
ruches et des métadonnées NTFS ; DPAPI, l'identité cloud, BITS, WMI, hosts,
setupapi, les volumes FAT, l'état de la machine en marche.

Sur les 1\,430 fichiers communs, 1\,412 ont la même empreinte SHA-1. Les 18
autres sont des fichiers que Windows réécrit sans cesse --- \$MFT, journal des
modifications, Prefetch, journaux de Defender, WebCache --- et qu'il a modifiés
entre les deux collectes.

\section{Ce qui est consigné de chaque pièce}

\begin{longtable}{@{}L{3.6cm}L{6.0cm}L{5.4cm}@{}}
\toprule
 & \textbf{\wac{} (manifeste)} & \textbf{KAPE (journal de copie)} \\
\midrule\endhead
Empreintes & MD5, SHA-1, SHA-256 & SHA-1 \\
Dates du fichier & les 4 de \texttt{\$STANDARD\_INFORMATION} et les 4 de
  \texttt{\$FILE\_NAME}, au 100\,ns, en UTC et en heure locale du suspect &
  3 de \texttt{\$STANDARD\_INFORMATION} (création, modification, accès), au
  100\,ns, en UTC, sans le changement de l'enregistrement MFT \\
Emplacement & numéro d'enregistrement MFT & chemin seul \\
Méthode & dite pour chaque pièce (lecture brute NTFS, observation vivante) &
  copie ordinaire ; lecture brute « différée » pour les fichiers verrouillés
  (180 sur 2\,674) \\
Échec & code et cause dans le manifeste, pièce par pièce & console et journal \\
Intégrité & manifeste scellé (SHA-256), revérifié à chaque conversion & --- \\
\bottomrule
\end{longtable}

\begin{note}
Les dates de \texttt{\$FILE\_NAME} sont celles qu'un intrus ne réécrit pas
avec les fonctions ordinaires : un fichier antidaté (« timestomping ») garde
dans \texttt{\$FILE\_NAME} sa vraie date de création. Le journal de KAPE ne les
a pas ; le manifeste de \wac{} les donne pour chaque pièce, avec l'enregistrement
MFT qui permet de la retrouver dans \fichier{mft.json}.
\end{note}

% ============================================================ 5
\chapter{Traces laissées sur la machine examinée}\label{chap:traces}

\section{La mesure}

Un clone de la machine de test, sans réseau, est figé avec Sysmon~15.22 installé ;
les outils sont sur une clé USB, lettre K:. Chaque exécution part de ce même
état : le disque système et la clé sont des couches neuves sur les images
figées. Le démarrage passé, chaque outil est lancé de la même façon --- une
tâche planifiée de l'utilisateur, avec ses droits élevés, comme un examinateur
depuis une invite élevée.

\begin{description}[style=nextline,leftmargin=1em]
  \item[Les fichiers] Le numéro USN du journal des modifications de C: est relevé
        juste avant et juste après l'outil : tout fichier créé, modifié, renommé
        ou supprimé entre les deux y figure, même créé puis supprimé pendant
        l'exécution. Les relevés sont écrits sur la clé, après la borne de fin.
  \item[L'auteur] Sysmon donne l'arbre des processus de l'outil et ce que chacun
        a écrit, fichiers et registre, et chaque processus qu'il a ouvert. Une
        écriture d'un autre processus qui nomme l'outil (Amcache, Prefetch, BAM)
        est une réaction de Windows à son lancement.
  \item[Le registre] Après l'arrêt propre de chaque exécution, les ruches du
        disque --- SYSTEM, SOFTWARE, SAM, SECURITY, DEFAULT, Amcache, celles des
        profils --- sont lues hors ligne, la machine éteinte, et confrontées
        clé par clé et valeur par valeur à celles de l'image figée. C'est le
        compte indépendant de ce que Sysmon attribue aux processus. Le volume
        système, chiffré par BitLocker sans protecteur, est ouvert par sa clé en
        clair.
  \item[Le témoin] Une exécution qui ne lance rien (\texttt{ver}) mesure le
        lanceur et l'activité propre de Windows ; ce qu'elle montre aussi n'est
        pas imputé à l'outil.
  \item[Le différé] Le ShimCache n'est écrit qu'à l'arrêt : chaque exécution se
        termine par un arrêt propre, un redémarrage et sa lecture.
\end{description}

\section{Ce que chacun laisse}

\begin{longtable}{@{}L{2.5cm}L{3.0cm}L{4.0cm}L{5.5cm}@{}}
\toprule
\textbf{Outil} & \textbf{Fichiers écrits sur C:} & \textbf{Registre écrit} &
  \textbf{Réactions de Windows} \\
\midrule\endhead
\wac{} (collecte) & aucun & aucun & Amcache : une entrée pendant
  l'exécution ; ShimCache : une entrée à l'arrêt \\
KAPE, cibles & aucun & 4 valeurs \texttt{ZoneMap} d'Internet Settings dans la
  ruche de l'utilisateur ; 115 ouvertures en création (magasins de
  certificats, stratégies) & Amcache et ShimCache pour \fichier{kape.exe} \\
KAPE, cibles et modules & 28 : 14 dans
  \fichier{ProgramData\textbackslash Microsoft\textbackslash NetFramework\textbackslash
  BreadcrumbStore}, 14 dans le Temp de l'utilisateur
  (\fichier{system-commandline-sentinel-files}) & les mêmes 4 valeurs,
  réécrites par 14 processus (136 écritures) ; 8 clés de traçage RAS créées
  par 4 outils (56 valeurs) ; 126 ouvertures en création & Amcache et
  ShimCache pour KAPE et 14 outils ; table des partitions de la clé
  réécrite dans USBSTOR \\
Outils EZ, en direct & aucun & aucune valeur ; 4 ouvertures en création
  (stratégies, \texttt{Tcpip}, \texttt{Internet Settings}) & Amcache et
  ShimCache pour les 6 outils \\
\bottomrule
\end{longtable}

Les ruches lues hors ligne le confirment. Hors ce que le témoin change aussi,
elles changent pendant chaque exécution --- de 59 clés et valeurs pour KAPE
à 311 pour \wac{}, dont la collecte dure 15 minutes ---, mais seules les
écritures que Sysmon attribue aux processus de l'outil lui reviennent : pour
\wac{}, aucune ; le reste est l'activité de Windows dans la fenêtre de mesure
(WaaSMedic, Defender, tâches planifiées, WMI, Windows Update, OneDrive,
l'Explorateur), sans lien avec l'outil, et l'entrée Amcache de chaque
exécutable lancé. Les clés de traçage RAS des modules de KAPE y apparaissent
ajoutées, comme Sysmon les donne.

Les processus des outils n'ont fait écrire aucun événement dans les journaux
de Windows eux-mêmes (System, Security, Application, journaux d'exploitation) :
ceux de la fenêtre de mesure sont les ouvertures de session des services,
Defender, la stratégie de groupe, le compte Microsoft, la télémétrie, communs
au témoin ou d'autres processus. L'audit de création de processus n'est pas
actif sur le clone ; actif, il aurait inscrit un 4688 et un 4689 par
processus (chapitre~\ref{chap:limites}).

La surveillance installée, elle, enregistre l'exécution de chaque outil : le
journal de Sysmon a reçu, pour les processus de l'outil seulement,

\begin{longtable}{@{}L{2.6cm}L{2.0cm}L{1.8cm}L{3.8cm}L{4.4cm}@{}}
\toprule
\textbf{Outil} & \textbf{Processus} & \textbf{Lectures brutes} & \textbf{Fichiers créés} & \textbf{Processus ouverts} \\
\midrule\endhead
\wac{} & 1 & 3 & 25\,532, tous sur la clé ; 777 supprimés sur la clé & 128 en lecture limitée (0x1000), 42 en requête et lecture mémoire (0x1410) \\
KAPE, cibles & 1 & 2 & 2\,940, tous sur la clé & Defender, en synchronisation \\
KAPE, cibles et modules & 15 & 2 & 3\,008 sur la clé, 28 sur C: & ses outils, en accès complet (17 ouvertures) ; Defender, en synchronisation (18) \\
Outils EZ, en direct & 6 & 6 & 45, tous sur la clé & Defender, en synchronisation (6) \\
\bottomrule
\end{longtable}

\noindent auxquels s'ajoutent les bibliothèques chargées par chaque processus
(25 pour \wac{}). Un EDR en tire la même télémétrie, et peut signaler la lecture
brute des volumes (\texttt{RawAccessRead}) : celles de \wac{} portent sur C:, la
partition EFI et le disque physique, pour les secteurs de démarrage.

\begin{description}[style=nextline,leftmargin=1em]
  \item[Processus ouverts] \wac{} lit la liste des processus et leurs
        propriétaires par Terminal Services, qui ouvre chacun en lecture
        limitée (0x1000), LSASS compris ; puis, pour la liste des modules de
        chaque processus, il l'ouvre en requête et lecture mémoire (0x1410),
        sans droit d'écriture, de fil ni d'injection. Les processus protégés
        (LSASS, services, Defender) refusent la seconde. KAPE ouvre ses
        propres outils en accès complet pour les piloter.
  \item[Ouverture en création] \texttt{RegCreateKeyEx} crée la clé si elle manque
        et l'ouvre sinon ; Sysmon ne distingue pas les deux cas. Ce sont les
        bibliothèques .NET qui ouvrent ainsi les magasins de certificats et les
        clés de stratégie. Une clé créée change de date ; une clé seulement
        ouverte, non.
  \item[Les valeurs \texttt{ZoneMap}] Écrites par l'initialisation de WinINet dans
        chaque processus .NET qui la charge, elles modifient la ruche
        \fichier{NTUSER.DAT} de l'utilisateur examiné, donc son journal et la date
        de la clé.
  \item[Les clés de traçage RAS] Le chargement des bibliothèques d'accès
        distant (\fichier{rasapi32.dll}) crée, pour chaque programme, deux clés
        sous \path|HKLM\SOFTWARE\Microsoft\Tracing| : EvtxECmd, MFTECmd, RECmd et
        SrumECmd lancés par KAPE en ont laissé huit dans SOFTWARE.
  \item[Les fichiers .NET] Le runtime laisse une « miette » par assemblage chargé
        dans \fichier{BreadcrumbStore}, et la bibliothèque System.CommandLine un
        fichier par outil dans le Temp de l'utilisateur. Lancés seuls, les
        outils EZ n'en ont pas laissé ; lancés par KAPE, si.
  \item[Prefetch et BAM] Aucun des programmes lancés depuis la clé USB n'a
        reçu de fichier Prefetch ni de valeur BAM, quand \fichier{cmd.exe}, lancé
        du disque système par la même tâche, a sa valeur BAM. Un outil copié sur la machine et lancé de
        là en recevrait.
  \item[Le prérequis .NET] KAPE et les outils EZ demandent le runtime .NET~9,
        présent sur le clone. Sur une machine qui ne l'a pas, l'installer écrit
        des centaines de fichiers et de clés ; \wac{} est un exécutable autonome.
\end{description}

\section{Ce que la mesure a corrigé dans \wac{}}

\wac{} effaçait l'écran au démarrage par \texttt{system("cls")}, qui lance
\fichier{cmd.exe} : un processus créé sur la machine examinée et le Prefetch de
\fichier{CMD.EXE} mis à jour, à chaque collecte. L'écran est désormais effacé par
l'API de la console, \wac{} n'importe plus aucune fonction capable de lancer un
processus, et le script de compilation refuse une version qui en importerait
une.

La mesure du 29 octobre 2026 a montré une seconde trace : pour lister les
modules de chaque processus, l'instantané de Toolhelp
(\texttt{CreateToolhelp32Snapshot}) ouvrait lui-même 42 processus de la
session en accès complet (\texttt{PROCESS\_ALL\_ACCESS}, 0x1FFFFF), le motif
même que surveille un EDR. \wac{} ouvre désormais chaque processus avec les
seuls droits de requête et de lecture mémoire, et lit ses modules par psapi ;
remesuré le 30 octobre, aucune ouverture en accès complet.

\begin{note}
Chaque trace de \wac{} est expliquée : il ne lance aucun processus, n'écrit
aucun fichier sur la machine, aucune valeur de registre, aucun événement dans
les journaux de Windows. Restent celles que Windows laisse de tout programme
lancé --- une entrée Amcache, une entrée ShimCache à l'arrêt ---, et ce que la
surveillance en place enregistre de son exécution (Sysmon, audit de création
de processus, EDR) : inévitables, et décrites dans le dossier d'innocuité.
\end{note}

% ============================================================ 6
\chapter{Précision}\label{chap:precision}

\section{En faveur de \wac{}}

\begin{description}[style=nextline,leftmargin=1em]
  \item[Journaux de ruche rejoués dans l'ordre] Les journaux \fichier{.LOG1} et
        \fichier{.LOG2} gardent, derrière les entrées récentes, les restes d'une
        génération précédente. RECmd les applique dans l'ordre du fichier, donc
        après les plus récentes. Il ramène ainsi des ruches à un état plus ancien :
        34 dernières exécutions BAM plus anciennes que celles de la ruche, la
        dernière connexion du compte de test et le dernier branchement des deux
        clés USB d'avant le démarrage du jour, des compteurs d'UserAssist en
        retard (mesure du 29 octobre 2026). \wac{} trie les entrées par séquence et
        s'arrête à la rupture de chaîne. Les journaux eux-mêmes le vérifient,
        sommes Marvin32 comprises.
  \item[Événements des journaux vidés] Un journal vidé garde ses enregistrements :
        \wac{} les relit et les marque \champ{EvtRecovered}, sans jamais les
        présenter comme vivants ; EvtxECmd ne les voit pas.
  \item[Valeurs d'événements comme Windows les rend] Hexadécimal \texttt{0x3e7},
        GUID entre accolades, booléens \texttt{true}, binaire \texttt{69060000},
        dates avec leur \texttt{Z}, marques U+200E conservées : les formes
        qu'affiche l'Observateur d'événements et que l'on recherche. EvtxECmd en
        change neuf, et se trompe sur 25 valeurs, par exemple un booléen
        \texttt{False} que Windows rend \texttt{true} ; chacune est tranchée par
        le XML que rend Windows (\texttt{Get-WinEvent}) sur le même journal.
  \item[Rien d'inventé] Une valeur absente de la source n'est pas émise.
        AmcacheParser écrit 5\,500 valeurs par défaut absentes de la ruche
        (\texttt{IsPeFile: False} pour un exécutable PE), SrumECmd 18\,412
        (0 ou $-1$).
  \item[SHA-1 partiel de l'Amcache signalé] Pour un fichier de plus de 30\,Mio,
        Amcache ne hache que les 31\,457\,280 premiers octets (mesuré sur
        \fichier{mrt.exe}). \wac{} le dit (\champ{AmcacheSha1CoversBytes}) ; sans cela,
        la comparaison avec le SHA-1 du fichier entier ferait conclure à tort à
        une modification.
  \item[Pas de champ trompeur] L'indicateur « exécuté » du ShimCache, que donne
        AppCompatCacheParser, dit « non » pour 105 des 117 binaires dont le
        Prefetch prouve l'exécution. \wac{} ne le relève pas : ce « non » ferait
        conclure qu'un programme n'a jamais tourné.
  \item[Précision au 100\,ns] \wac{} donne chaque date à la précision de sa
        source. Par défaut, PECmd, LECmd et SrumECmd arrondissent à la seconde, et
        LECmd ignore \opt{mp} dans son CSV.
  \item[Heure locale datée] Chaque date est donnée en UTC et en heure locale du
        suspect, avec le décalage en vigueur à cette date. Les outils EZ ne
        donnent que l'UTC, sans le signaler par un \texttt{Z}.
  \item[Lectures d'autrui vérifiées] Les « FirstInteracted » et
        « LastInteracted » de SBECmd, le « CreatedOn » des comptes chez RECmd sont
        des lectures de dates d'écriture de clés. \wac{} publie ces dates
        elles-mêmes : chacun en tire la lecture qu'il veut, sans en hériter une
        toute faite. SBECmd donne en outre un \texttt{NodeSlot} faux (0) aux
        dossiers de premier niveau.
  \item[Erreurs de lecture ESE de SrumECmd] 8 lignes d'énergie sur 2\,306 portent
        le contenu de la ligne précédente ; dissect.esedb donne raison à \wac{}.
  \item[Base ESE non fermée proprement] Une base SRUM prise sur une machine en
        marche peut ne pas être fermée proprement. SrumECmd s'est arrêté sur
        l'une d'elles (mesure du 9 octobre 2026) ; \wac{}
        rejoue ses journaux sur la copie de travail (\texttt{esent.dll}), puis la
        lit.
  \item[Comptes supprimés] Le banc supprime un compte (\texttt{wacrdp}) après
        s'en être servi. RECmd le liste encore, d'une SAM ramenée à un état
        antérieur ; \wac{} ne le donne pas, et \texttt{Get-LocalUser} non plus.
  \item[Collecte] Chaque pièce de \wac{} porte MD5, SHA-1 et SHA-256, les huit
        dates NTFS de son original et son entrée MFT, dans un manifeste scellé.
        KAPE donne un SHA-1 et trois dates.
\end{description}

\section{En faveur des outils EZ}\label{sec:couverture}

\begin{description}[style=nextline,leftmargin=1em]
  \item[Couverture] Il n'y a plus de lacune ouverte. Les clés du lot Kroll que
        \wac{} ignorait sont désormais lues : programmes installés,
        exclusions de Defender, RecentDocs, MountPoints2, réseaux connus,
        configuration des journaux, entre autres. Celles qui ne tracent aucune
        activité sont écartées, avec leur raison.
        Les fichiers que KapeTriage était seul à prendre sont collectés et
        convertis : navigateurs Chromium, WebCache et INetCache, sessions et
        onglets ouverts ou fermés, favicons, prédicteur de la barre d'adresse,
        chronologie Windows, raccourcis hors
        de Recent, la ruche du profil Default, \fichier{\$Secure:\$SDS},
        \fichier{\$Boot}, les traces WPP de Defender et les clés maîtresses
        DPAPI. Seuls les \fichier{desktop.ini} sont écartés, faute de valeur
        d'enquête.
  \item[Vitesse sur les petits artefacts] Pris un par un, Prefetch, Amcache et
        raccourcis se décodent plus vite avec l'outil EZ : voir le
        chapitre~\ref{chap:temps}.
  \item[Interprétations prêtes à lire] Les outils EZ découpent et nomment certains
        champs : exécutable, date et description d'une application SRUM, noms des
        dossiers connus, extension d'un fichier. \wac{} donne la valeur brute,
        telle que Windows la stocke.
\end{description}

\begin{attention}
Ce que KAPE prend sans que \wac{} le prenne, ce sont surtout des fichiers
d'applications dans les profils (chapitre~\ref{chap:prelevement}). Ils ne
correspondent à aucun artefact que \wac{} décode : leur intérêt dépend de
l'enquête.
\end{attention}

% ============================================================ 7
\chapter{Formats de sortie}

\section{Ce que produit chaque outil}

\begin{longtable}{@{}L{4.4cm}L{5.4cm}L{5.2cm}@{}}
\toprule
\textbf{Outil} & \textbf{Formats} & \textbf{Remarque} \\
\midrule\endhead
AmcacheParser & CSV & 8 fichiers, un par table \\
AppCompatCacheParser & CSV & \\
RBCmd & CSV & \\
SrumECmd & CSV & 7 fichiers \\
EvtxECmd & CSV, JSON, XML & \\
JLECmd & CSV, JSON, HTML & \\
LECmd & CSV, JSON, HTML, XML & \\
MFTECmd & CSV, JSON, bodyfile & \\
PECmd & CSV, JSON, HTML & \\
RECmd & CSV, JSON & plus un CSV par greffon \\
SBECmd & CSV, JSON & \\
KAPE & copie des fichiers, journal de copie CSV & \\
\wac{} & JSON pour chaque artefact & un fichier par artefact, plus le manifeste scellé \\
\bottomrule
\end{longtable}

\section{Ce qui les distingue}

Quatre outils sur onze ne produisent que du CSV. Chez les autres, le JSON est la
transcription de la ligne CSV : une ligne de texte par enregistrement (JSON
Lines), à plat.

\begin{description}[style=nextline,leftmargin=1em]
  \item[Types] Chez les outils EZ, les nombres sont des chaînes : un compteur
        d'exécutions s'écrit \texttt{"RunCount":"458"}, un numéro
        d'enregistrement \texttt{"EventRecordId":"21572"}. \wac{} écrit des
        nombres, des booléens, des chaînes.
  \item[Structures] Chez les outils EZ, une liste devient une chaîne jointe par
        des virgules (les fichiers chargés d'un Prefetch, ses répertoires), et les
        données d'un événement sont un JSON échappé \emph{dans} une chaîne
        (\texttt{Payload}), à décoder une seconde fois. \wac{} imbrique : les
        volumes d'un Prefetch et leurs répertoires, les éléments d'un raccourci,
        les paires \champ{Name}/\champ{Value} d'un événement.
  \item[Dates] Chez les outils EZ, l'UTC n'est pas signalé
        (\texttt{2026-09-30 23:04:46}) et la précision dépend d'une option. \wac{}
        donne l'ISO 8601 complet, en UTC (\texttt{…Z}) et en heure locale avec son
        décalage (\texttt{…+02:00}).
  \item[Absence de valeur] Les outils EZ écrivent une valeur par défaut, ou une
        cellule vide. \wac{} n'émet pas le champ ; un artefact non collecté est
        déclaré tel (\texttt{NotRequested}, \texttt{NotCollected}), pour qu'une
        liste vide ne passe pas pour une absence de trace.
  \item[Provenance] Le manifeste de \wac{} relie chaque valeur à la pièce dont
        elle vient, avec ses empreintes et ses dates d'origine.
  \item[Console] Les outils EZ détaillent leur travail sur la console ; certains
        y écrivent le décodage lui-même quand aucun fichier n'est demandé. \wac{}
        n'y écrit que l'avancement : tout ce qu'il trouve est dans les fichiers.
\end{description}

\begin{note}
Pour un analyste, la différence se voit au premier traitement : un JSON de
\wac{} se charge tel quel dans un outil d'analyse ou dans WAE, sans rien
redécouper ; une sortie EZ se traite comme un tableur, dont il faut reconvertir
les nombres, redécouper les listes et redécoder les données d'événement.
\end{note}

% ============================================================ 8
\chapter{Ce que la confrontation a appris à \wac{}}\label{chap:defauts}

Chaque défaut ci-dessous produisait un JSON valide et faux, ou incomplet. Tous
sont corrigés, avec un contrôle qui les aurait pris.

\begin{longtable}{@{}L{3.0cm}L{12.0cm}@{}}
\toprule
\textbf{Artefact} & \textbf{Défaut corrigé} \\
\midrule\endhead
\$MFT & enregistrements d'extension périmés fusionnés avec la base \\
EVTX & flottants tronqués à 6 chiffres ; binaire en « 69 06 00 00 » ; entiers hexadécimaux et GUID au mauvais format ; ProcessingErrorData ignoré (21 événements sans données) ; événements des journaux vidés non lus \\
Raccourcis & URI collée au bloc qui la suit ; lecture décalée de 2 octets sans liste d'ID ; blocs ExtraData non lus (machine, suivi) ; taille de la cible non lue ; DROID nuls émis \\
Jump lists & nombre d'ouvertures et rang MRU non lus ; format de DestList version 1 non pris en charge ; nom d'application altéré \\
Shellbags & dossier perdu avec sa descendance ; date de la clé du dossier, rang et NodeSlot absents ; dossier « Accueil » sans nom ; arbre de \fichier{NTUSER.DAT} non lu (seul celui de \fichier{UsrClass.dat} l'était) \\
Amcache & 2 tables sur 14 et 4 à 6 valeurs lues : ni SHA-1, ni taille, ni pilote, ni périphérique ; SHA-1 nommé \champ{Sha1}, lu comme celui du fichier entier, devenu \champ{AmcacheSha1} \\
ShimCache & tabulations des applications empaquetées remplacées par des espaces \\
SRUM & dates FILETIME laissées en nombres ; texte UTF-16 lu comme de l'ANSI \\
Identifiants connus & 105 chemins de DLL de la table des GUID (253 occurrences) dont la génération avait pris \texttt{\textbackslash n} ou \texttt{\textbackslash r} pour un saut de ligne : \fichier{system32\textbackslash netcorehc.dll} devenu \fichier{system32 etcorehc.dll} \\
Registre & BAM limité aux comptes ayant un profil (un compte supprimé perdait ses exécutions) ; durée de premier plan de UserAssist non lue \\
Collecte & pièces dont le chemin dépasse 260 caractères une fois dans la consigne non écrites (482 sur le kit complet) : écrites désormais en chemin étendu \\
Lancement & \texttt{system("cls")} lançait \fichier{cmd.exe} sur la machine examinée (mesure des traces) \\
Processus & la liste des modules de chaque processus, par l'instantané de Toolhelp, faisait ouvrir 42 processus en accès complet (mesure des traces) : ouverts désormais en requête et lecture mémoire seules \\
Harnais & un contrôle qui plantait interrompait tous les suivants sans échec visible ; les outils EZ lancés sur le magasin scellé lui-même --- SrumECmd y a écrit, \wac{} a refusé la conversion suivante ; ils lisent désormais une copie ; LECmd confronté aux seuls raccourcis de Recent, la table PushNotifications de SRUM jamais appariée \\
\bottomrule
\end{longtable}

% ============================================================ 9
\chapter{Limites de la comparaison}\label{chap:limites}

\begin{itemize}
  \item Une seule machine, virtuelle, sous Windows 11 : certains artefacts y
        sont peu nombreux (2 clés USB, une recherche dans l'Explorateur). Les formats sont
        les mêmes sur une machine réelle, mais les cas rares n'y sont pas tous
        représentés.
  \item Les deux collectes se suivent au lieu d'être simultanées : les fichiers
        écrits entre elles (\$MFT, journaux, Prefetch) diffèrent légitimement.
  \item Les outils EZ ont été lancés avec leurs options de précision maximale ;
        par défaut, ils perdent davantage.
  \item Les temps viennent d'une machine virtuelle et de son disque : ils
        comparent les outils entre eux, ils ne prédisent pas la durée sur une
        machine réelle, où le nombre d'exécutables et la vitesse du disque
        pèsent le plus.
  \item La mesure des traces se fait sur un clone sans réseau, en une
        exécution par outil : une écriture de Windows que le témoin n'a pas
        montrée n'est pas imputée à l'outil sans que Sysmon l'attribue à l'un de
        ses processus ; elle reste listée comme non attribuée dans les résultats
        de \fichier{footprint-diff.py}.
  \item L'audit de création de processus (Security 4688) n'est pas actif sur la
        machine de test : sur une machine où il l'est, chaque outil y laisse un
        événement par processus --- un pour \wac{}, quinze pour KAPE et ses
        modules.
  \item La comparaison porte sur les artefacts communs. Ce que seul l'un des
        outils lit n'a pas de référence ici : son exactitude repose sur les
        contrôles du dossier de validation.
\end{itemize}

\end{document}
